\begin{abstract}
Training a 2B model to judge its own mathematical attempts raised its AUROC from 0.57 to 0.67, yet picking among its
answers did not improve. We show this is predicted, not anomalous. Three measurement illusions hide what a
self-verifier can do:
(i) a lenient judge looks near chance under sampled verdicts even when its underlying score discriminates;
(ii) on a typical pool, 99.9\% of the pairs behind an AUROC compare different questions, while selection and repair
provably use only within-question information; and
(iii) majority vote saturates, and before that point verification cannot beat it at equal compute.

A five-parameter model fitted only to 4-attempt aggregates predicts the 16-attempt scaling curves of 16,000 new
attempts to within a point. Reading verdict probabilities instead of sampling verdicts raises the judge's AUROC by 0.07
to 0.09 (paired 95\% CIs exclude zero). Judge-free trace features match the trained judge's pooled AUROC yet are at
chance on contested questions, exactly as the theory requires.

The theory predicts a procedure, contest-adaptive verification, that spends judgement only where sampled answers
disagree. In a pre-registered test on 800 questions it beats majority vote at equal compute by $+1.6$ points
$[+0.5, +2.6]$ and passes the plateau that voting and adaptive-consistency share, close to the $\approx+1.3$ the
model predicted before judging.
\end{abstract}
